\documentclass[10pt]{article}
\usepackage{newtxtext}
\usepackage{amsmath}
\usepackage{geometry}
\geometry{a4paper, left=0.75in, right=0.75in, top=0.55in, bottom=0.55in}
\usepackage{graphicx}
\usepackage{parskip}
\usepackage{enumitem}
\usepackage[colorlinks=true, urlcolor=blue, linkcolor=blue, citecolor=blue]{hyperref}
\usepackage{xurl}

\begin{document}
\pagestyle{empty}

% --- Header / Letterhead ---
\begin{minipage}[b]{0.32\textwidth}
  \includegraphics[height=1.8cm]{logo.jpeg}
\end{minipage}%
\begin{minipage}[b]{0.68\textwidth}
  \raggedleft
  {\large\bfseries Jason Pandian} \\
  Department of Information Technology \\
  Nehru Institute of Technology \\
  Coimbatore, Tamil Nadu, 641105, India \\
  \textbf{Email: pandiajason@gmail.com}
\end{minipage}

\vspace{0.05cm}
\hrule height 1pt
\vspace{0.08cm}

\hfill \begin{tabular}{ll}
  \textbf{Date:} & October 07, 2026 \\
  \textbf{Place:} & Coimbatore, India
\end{tabular}

\vspace{0.15cm}

\noindent
\textbf{To:} \\
The Editor-in-Chief and Area Editor \\
\textit{Simulation Modelling Practice and Theory} \\
Elsevier Science

\vspace{0.15cm}

\noindent
\textbf{Subject: Submission of Revised Manuscript [Ref. No.: SIMPAT-D-26-1229]}

\vspace{0.12cm}

\noindent
Dear Editor-in-Chief and Area Editor,

Thank you for the opportunity to revise our manuscript, \textbf{SIMPAT-D-26-1229}, titled \textbf{``The Limits of Discrete-Event Synchronization in Interplanetary Network Simulation: A Critical Review of Co-Simulation Fidelity,''} authored by \textbf{Jason Pandian} and \textbf{I. Kala}.

We express our sincere gratitude to the Area Editor for guiding our paper and to the two anonymous reviewers for their thorough, rigorous, and constructive evaluations. Their insightful recommendations have enabled us to substantially strengthen the theoretical framework, algorithmic precision, empirical benchmarks, and scholarly grounding of our work.

In accordance with your editorial decision and instructions, we have performed a comprehensive revision. We have compiled a detailed, point-by-point \textbf{Response to Reviewers} addressing all comments, and we have \textbf{explicitly highlighted all modifications in the revised manuscript} using blue text to facilitate editorial re-evaluation.

\vspace{0.08cm}
\noindent
\textbf{\large Executive Summary of Major Revisions}
\vspace{0.05cm}

\begin{enumerate}[leftmargin=1.8em, itemsep=2.5pt, topsep=1pt, parsep=0pt]
  \item \textbf{Correction and Attribution of Speedup Metrics (Reviewer 1, Comment 1):}
  We have corrected all reporting across the Abstract, Highlights, Table 3, and Sections 5.3--5.4. The peak speedup achieved by our proposed adaptive engine (\texttt{CoSimOrchestrator}) is accurately stated as \textbf{$9.90\times$} on Cislunar ($12.3\times$ reduction in integration steps), \textbf{$8.05\times$} on the Elliptical relay, and \textbf{$5.43\times$} on Mars. The extreme $435.4\times$ speedup on Cislunar is strictly and unambiguously attributed to the zero-order hold Event-Driven baseline. Furthermore, we explicitly emphasize why the baseline's apparent speedup is physically unviable: freezing physics updates during communication blackouts causes $>30,970$~km of uncorrected orbital drift, introducing $6.2$--$14.1$~dB link margin errors and false ModCod rate selections, whereas our adaptive method preserves meter-scale physical accuracy ($32.68$~m).

  \item \textbf{Separation of Relativistic Physics Limits from Runtime Execution (Reviewer 1, Comment 2):}
  To clearly distinguish between implemented algorithmic routines and theoretical problem bounds, we removed Section 3.5 from System Modeling. Equations 1--3 detailing relativistic proper-time drift (Fidelity Gap 3) have been relocated to the Introduction (Section 1.1) to define the theoretical limits of coordinate-time simulation. An explicit scope note clarifies that our prototype orchestrator couples in Barycentric Dynamical Time (TDB), while Equations 1--3 serve as the theoretical ceiling and mathematical blueprint for future proper-time-aware architectures.

  \item \textbf{Algorithmic Boundary Event Guard and Mars Speedup Resolution (Reviewer 1, Comment 3):}
  Addressing the bottleneck observed on the Mars profile caused by step-size overshoot and boundary chattering near ingress/egress, we formulated and implemented a predictive boundary event guard rule:
  \begin{equation*}
    \Delta t = \min(\Delta t_{\text{proposed}},\, t_{\text{next\_event}} - t)
  \end{equation*}
  This rule proactively snaps the simulation step to the exact event boundary. Consequently, Mars integration steps plummeted from 486 to just \textbf{49 steps}, execution time dropped to $0.0169$~s (yielding a \textbf{$5.43\times$} speedup over fixed-step), and occultation boundary timestamp error was completely eliminated ($0.0$~ms). An analytical boundary analysis was added in Section 5.3.2.

  \item \textbf{Validation Under Extreme Network Burst Loads and Micro-Events (Reviewer 1, Comment 4):}
  We clarified the decouple-and-intercept mechanism in Sections 3.4 and 4.3. Under line-of-sight packet bursts, discrete queue drops and custody releases execute discretely within the network engine without corrupting the macro orbital step. Furthermore, any scheduled micro-event occurring prior to $\Delta t$ contracts the coupling step, ensuring zero skipped micro-events, validated against a fine-step reference ($\Delta t_{\text{ref}} = 0.01$~s).

  \item \textbf{Discrete CCSDS Variable Coding and Modulation (ModCod) Model (Reviewer 1, Comment 5):}
  We incorporated a discrete ModCod staircase mapping table (QPSK, 8PSK, 16APSK, 32APSK with LDPC codes) compliant with CCSDS 131.2-B-1 and 231.0-B-1 (Section 3.2, Equation 13). Continuous Shannon capacity is retained solely as an analytical reference upper bound (Equation 14). We show that discrete ModCod transitions necessitate adaptive stepping, as fixed steps miss modulation switching thresholds.

  \item \textbf{Multi-Node Constellation Scalability from 10 to 500 Spacecraft (Reviewer 1, Comment 6):}
  We expanded our empirical scalability benchmark by more than an order of magnitude, evaluating constellation sizes from 10 to 500 spacecraft nodes over the full 13-day Cislunar NRHO mission arc (Section 5.3.1, Figure 4b). Wall-clock runtime scaled sub-linearly from $0.789$~s for 10 nodes to $6.353$~s for 500 nodes (a $50\times$ node expansion required only an $8.0\times$ runtime increase), confirming that discrete queue operations do not degrade continuous orbital propagation.

  \item \textbf{Discrete Packet Queue Formulation for DTN Buffering (Reviewer 1, Comment 7):}
  We replaced the continuous fluid approximation with a discrete packet queue model (Section 3.3, Equation 16), summing over queued packets with explicit 1024-byte MTU boundaries, integer byte allocations, and discrete custody transfer timers. The scope of our model is clearly delineated as a temporal queue synchronization framework rather than a full-stack RFC 5050/9171 protocol emulator.

  \item \textbf{Substantial Modernization and Literature Expansion (Reviewer 2, Comment 1):}
  We expanded the bibliography from 22 to \textbf{39 references} (a $77\%$ increase), incorporating authoritative standards and peer-reviewed studies published between 2017 and 2024 (with 14 citations from 2020--2024). Newly integrated works include FMI 3.0 multi-rate clock specifications, 6G satellite digital twins (IEEE Network 2024), Hypatia, RFC 9171 BPv7, CCSDS 734.20-B-1, NASA LunaNet, and the White House OSTP Coordinated Lunar Time directive (April 2024).

  \item \textbf{Explicit Delineation of Original Contributions (Reviewer 2, Comment 2):}
  We introduced Section 1.2 (``Scope, Original Contributions, and Relation to Prior Work''), explicitly positioning our contributions against foundational co-simulation surveys (Gomes et al. [19], Schweiger et al. [20]) across four key axes: continuous-to-discrete cross-domain coupling, predictive look-ahead vs. reactive rollback, domain-specific multi-rate stepping regimes with event guards, and relativistic spacetime limits.

  \item \textbf{Figure Typography and Layout Standardization (Reviewer 2, Comment 3):}
  We removed all embedded canvas supertitles (\texttt{fig.suptitle()}) from the plotting scripts, regenerated all vector figures (Figures 1--6) with clean canvas margins, and ensured all captions are typeset strictly below each float per journal styling.
\end{enumerate}

\vspace{0.1cm}
\noindent
\textbf{\large Submission Package}
\vspace{0.05cm}

In this revised submission package, we provide:
\begin{itemize}[leftmargin=1.8em, noitemsep, topsep=1pt]
  \item \textbf{Revised Manuscript (Clean Version):} Formatted using the Elsevier \texttt{cas-sc.cls} template.
  \item \textbf{Revised Manuscript (Marked-Up Version):} Highlights all additions and revisions in blue text for convenient editorial review.
  \item \textbf{Detailed Response to Reviewers:} Complete point-by-point rebuttal with verbatim reviewer comments, author responses, and line-by-line change listings.
  \item \textbf{Open-Source Reproduction Package:} Complete Python codebase, SPICE kernel integration scripts, and benchmark datasets available at: \\
  \url{https://github.com/PandiaJason/SPICE-ns-Project/tree/main/Adaptive-Multi-Rate-Co-Simulation-Engine}.
\end{itemize}

\vspace{0.08cm}
\noindent
\textbf{\large Confirmations and Declarations}
\vspace{0.05cm}

We confirm that:
\begin{itemize}[leftmargin=1.8em, noitemsep, topsep=1pt]
  \item The revised manuscript has been thoroughly reviewed and approved by all authors.
  \item This manuscript represents original research and is not under consideration elsewhere.
  \item The authors have no competing commercial or financial interests to declare.
\end{itemize}

We trust that the comprehensive revisions and clarifications fully address the reviewers' comments and meet the high standards of \textit{Simulation Modelling Practice and Theory}. Thank you for your continued editorial guidance.

\vspace{0.2cm}
\noindent
Sincerely,

\vspace{0.15cm}
\begin{minipage}[b]{0.48\textwidth}
  \includegraphics[height=1.0cm]{Jason_Signature.png}
\end{minipage}\hfill
\begin{minipage}[b]{0.48\textwidth}
  \vspace{1.0cm}
\end{minipage}

\vspace{0.1cm}

\begin{minipage}[t]{0.48\textwidth}
  \textbf{Jason Pandian} (Corresponding Author) \\
  Assistant Professor / Researcher \\
  Department of Information Technology \\
  Nehru Institute of Technology \\
  Email: \texttt{pandiajason@gmail.com}
\end{minipage}\hfill
\begin{minipage}[t]{0.48\textwidth}
  \textbf{Dr. I. Kala} (Co-Author) \\
  Associate Professor \\
  Department of Computer Science and Engineering \\
  PSG Institute of Technology and Applied Research \\
  Email: \texttt{kala@psgitech.ac.in}
\end{minipage}

\end{document}
